\documentclass[11pt]{article}

\usepackage[margin=1in]{geometry}
\usepackage{amsmath, amssymb}
\usepackage{booktabs}
\usepackage{graphicx}
\usepackage{float}
\usepackage{hyperref}

\graphicspath{{outputs/figures/}}
\setlength{\parskip}{0.6em}
\setlength{\parindent}{0pt}

\title{Math 104C Project 2: Numerical Methods for Partial Differential Equations}
\author{Jiajia Liu \\ Perm ID: A0J2E85}
\date{}

\begin{document}
\maketitle

\textbf{Code link:} \href{https://github.com/samanthaliujiajia/math104c-project2}{https://github.com/samanthaliujiajia/math104c-project2}

\section{Introduction}

This project investigates finite-difference methods for elliptic and parabolic
partial differential equations. The experiments compare accuracy, stability,
mesh refinement, and computational behavior for the five-point finite-difference
method, Forward Difference, Backward Difference, and Crank--Nicolson.

The elliptic experiments use square grids with \(h=1/4\), \(h=1/8\), and
\(h=1/16\). The heat equation experiments use multiple choices of \(h\) and
\(k\), including both stable and unstable Forward-Difference cases. Whenever an
exact solution is available, the absolute error is computed at mesh points.

\section{Numerical Methods}

For elliptic equations on the unit square, the five-point stencil approximates
\[
  u_{xx}(x_i,y_j)+u_{yy}(x_i,y_j)
  \approx
  \frac{w_{i+1,j}+w_{i-1,j}+w_{i,j+1}+w_{i,j-1}-4w_{i,j}}{h^2}.
\]
Boundary values enter the right-hand side of the resulting sparse linear
system. The experiments solve the system both by a sparse direct solver and by
successive over-relaxation (SOR).

For the heat equation \(u_t=u_{xx}\), the Forward-Difference method is
\[
  w_{i,j+1}=\lambda w_{i-1,j}+(1-2\lambda)w_{i,j}+\lambda w_{i+1,j},
  \qquad \lambda=\frac{k}{h^2}.
\]
It is conditionally stable for \(\lambda\leq 1/2\). The Backward-Difference
method and Crank--Nicolson method require solving tridiagonal linear systems at
each time step, but they remain stable for much larger time steps.

\section{Problem A: Elliptic PDE}

Problem A solves
\[
  u_{xx}+u_{yy}=0,\qquad 0<x<1,\quad 0<y<1,
\]
with boundary conditions
\[
  u(0,y)=0,\quad u(x,0)=0,\quad u(x,1)=100x,\quad u(1,y)=100y.
\]
The exact solution is \(u(x,y)=100xy\), which is bilinear. Since the five-point
stencil is exact for this function up to roundoff, this problem is a useful
check that boundary values and matrix assembly are correct.

\begin{table}[H]
\centering
\caption{Elliptic PDE error summary and SOR iteration counts.}
\resizebox{\textwidth}{!}{\begin{tabular}{llllll}
\toprule
problem & method & $h$ & max error & iterations & residual \\
\midrule
A & Direct & 0.250000 & 7.105e-15 & 0 & 2.842e-14 \\
A & SOR omega=1.6 & 0.250000 & 5.321e-11 & 56 & 1.873e-10 \\
A & Direct & 0.125000 & 2.132e-14 & 0 & 3.553e-14 \\
A & SOR omega=1.6 & 0.125000 & 6.425e-11 & 58 & 1.914e-10 \\
A & Direct & 0.062500 & 2.132e-14 & 0 & 8.527e-14 \\
A & SOR omega=1.6 & 0.062500 & 8.285e-10 & 137 & 9.442e-11 \\
C1 & Direct & 0.250000 & 0.053029 & 0 & 8.882e-16 \\
C1 & SOR omega=1.6 & 0.250000 & 0.053029 & 47 & 1.591e-10 \\
C1 & Direct & 0.125000 & 0.012951 & 0 & 6.384e-16 \\
C1 & SOR omega=1.6 & 0.125000 & 0.012951 & 48 & 8.876e-11 \\
C1 & Direct & 0.062500 & 0.003219 & 0 & 1.527e-15 \\
C1 & SOR omega=1.6 & 0.062500 & 0.003219 & 114 & 8.457e-11 \\
\bottomrule
\end{tabular}
}
\end{table}

\begin{table}[H]
\centering
\caption{Observed orders for elliptic experiments.}
\begin{tabular}{lll}
\toprule
problem & method & observed order \\
\midrule
A & Direct & -0.792481 \\
A & SOR omega=1.6 & -1.980335 \\
C1 & Direct & 2.021060 \\
C1 & SOR omega=1.6 & 2.021061 \\
\bottomrule
\end{tabular}

\end{table}

\begin{figure}[H]
\centering
\includegraphics[width=\textwidth]{problem_A_Direct_h0p125.pdf}
\caption{Problem A exact solution, numerical solution, and absolute error for \(h=1/8\).}
\end{figure}

\section{Problem B: Heat Equation}

Problem B solves
\[
  u_t=u_{xx},\qquad 0<x<1,\quad t>0,
\]
with \(u(0,t)=u(1,t)=0\) and \(u(x,0)=\sin(\pi x)\). The exact solution is
\[
  u(x,t)=e^{-\pi^2t}\sin(\pi x).
\]
The Forward-Difference experiments include \(\lambda=0.4\), which satisfies the
stability condition, and \(\lambda=0.6\), which violates it.

\begin{table}[H]
\centering
\caption{Heat equation error summary.}
\resizebox{\textwidth}{!}{\begin{tabular}{llllllll}
\toprule
problem & case & method & $h$ & $k$ & $\lambda$ & final error & max error \\
\midrule
B & B\_stable\_h0.1 & Forward Difference & 0.100000 & 0.004000 & 0.400000 & 0.004294 & 0.004294 \\
B & B\_stable\_h0.1 & Backward Difference & 0.100000 & 0.004000 & 0.400000 & 0.010112 & 0.010112 \\
B & B\_stable\_h0.1 & Crank-Nicolson & 0.100000 & 0.004000 & 0.400000 & 0.002981 & 0.002981 \\
B & B\_refined\_h0.05 & Forward Difference & 0.050000 & 0.001000 & 0.400000 & 0.001063 & 0.001063 \\
B & B\_refined\_h0.05 & Backward Difference & 0.050000 & 0.001000 & 0.400000 & 0.002561 & 0.002561 \\
B & B\_refined\_h0.05 & Crank-Nicolson & 0.050000 & 0.001000 & 0.400000 & 7.535e-04 & 7.535e-04 \\
B & B\_unstable\_h0.1 & Forward Difference & 0.100000 & 0.010000 & 1.000000 & 0.015756 & 0.015756 \\
B & B\_unstable\_h0.1 & Backward Difference & 0.100000 & 0.010000 & 1.000000 & 0.020320 & 0.020320 \\
B & B\_unstable\_h0.1 & Crank-Nicolson & 0.100000 & 0.010000 & 1.000000 & 0.002734 & 0.002734 \\
C2 & C2\_stable\_h0.1 & Forward Difference & 0.100000 & 0.004000 & 0.400000 & 0.003285 & 0.017047 \\
C2 & C2\_stable\_h0.1 & Backward Difference & 0.100000 & 0.004000 & 0.400000 & 0.008843 & 0.036656 \\
C2 & C2\_stable\_h0.1 & Crank-Nicolson & 0.100000 & 0.004000 & 0.400000 & 0.002355 & 0.010840 \\
C2 & C2\_refined\_h0.05 & Forward Difference & 0.050000 & 0.001000 & 0.400000 & 8.740e-04 & 0.004294 \\
C2 & C2\_refined\_h0.05 & Backward Difference & 0.050000 & 0.001000 & 0.400000 & 0.002181 & 0.010114 \\
C2 & C2\_refined\_h0.05 & Crank-Nicolson & 0.050000 & 0.001000 & 0.400000 & 6.247e-04 & 0.002981 \\
\bottomrule
\end{tabular}
}
\end{table}

\begin{table}[H]
\centering
\caption{Observed orders for stable heat equation refinement cases.}
\begin{tabular}{lll}
\toprule
problem & method & observed order \\
\midrule
B & Backward Difference & 1.981501 \\
B & Crank-Nicolson & 1.983931 \\
B & Forward Difference & 2.014890 \\
C2 & Backward Difference & 1.857639 \\
C2 & Crank-Nicolson & 1.862650 \\
C2 & Forward Difference & 1.989060 \\
\bottomrule
\end{tabular}

\end{table}

\begin{figure}[H]
\centering
\begin{minipage}{0.48\textwidth}
\includegraphics[width=\textwidth]{B_stable_h0.1_final_profiles.pdf}
\centering Stable Forward case, \(\lambda=0.4\).
\end{minipage}
\begin{minipage}{0.48\textwidth}
\includegraphics[width=\textwidth]{B_unstable_h0.1_final_profiles.pdf}
\centering Unstable Forward case, \(\lambda=0.6\).
\end{minipage}
\caption{Problem B final-time solution profiles.}
\end{figure}

\begin{figure}[H]
\centering
\includegraphics[width=0.74\textwidth]{problem_B_convergence.pdf}
\caption{Problem B max error under stable mesh refinement.}
\end{figure}

\section{Problem C: Additional Experiments}

The first additional experiment solves a Poisson equation with a manufactured
solution:
\[
  u_{xx}+u_{yy}=-2\pi^2\sin(\pi x)\sin(\pi y),\qquad
  u=0 \text{ on the boundary}.
\]
The exact solution is \(u(x,y)=\sin(\pi x)\sin(\pi y)\). Unlike Problem A, this
solution is not bilinear, so the second-order spatial truncation error is
visible in the grid refinement results.

\begin{figure}[H]
\centering
\includegraphics[width=\textwidth]{problem_C1_Direct_h0p125.pdf}
\caption{Poisson manufactured-solution experiment for \(h=1/8\).}
\end{figure}

\begin{figure}[H]
\centering
\includegraphics[width=0.74\textwidth]{problem_C1_convergence.pdf}
\caption{Problem C1 max error under mesh refinement.}
\end{figure}

The second additional experiment changes the heat equation initial condition to
\[
  u(x,0)=\sin(2\pi x),
\]
whose exact solution is
\[
  u(x,t)=e^{-4\pi^2t}\sin(2\pi x).
\]
This higher-frequency mode decays faster and is more demanding for a coarse
mesh.

\begin{figure}[H]
\centering
\includegraphics[width=0.74\textwidth]{C2_stable_h0.1_final_profiles.pdf}
\caption{Problem C2 final-time profiles for the different initial condition.}
\end{figure}

\begin{figure}[H]
\centering
\includegraphics[width=0.74\textwidth]{problem_C2_convergence.pdf}
\caption{Problem C2 max error under stable mesh refinement.}
\end{figure}

\section{Discussion and Analysis}

Problem A confirms both the finite-difference stencil and the treatment of
boundary data. Since the exact solution \(u(x,y)=100xy\) is bilinear, the
second derivatives are represented exactly by the centered difference formula up
to roundoff error. Therefore, the small errors in the table are not surprising:
they mainly measure floating-point error and the accuracy of the linear solve,
not truncation error. This is still an important test because an incorrect sign
in the matrix or a missing boundary contribution would immediately destroy the
linear solution.

The direct solver and SOR solver give essentially the same numerical solution,
but their computational behavior is different. The sparse direct solver treats
the linear system as one algebraic problem. SOR updates the grid values
repeatedly, so the number of iterations grows as the mesh is refined. This is a
typical feature of stationary iterative methods: finer grids have more unknowns,
and information moves across the domain gradually through local updates. The
residuals show that the iteration was stopped only after the algebraic system
was solved to a tight tolerance.

Problem B illustrates the central stability difference between explicit and
implicit time stepping. For the heat equation, the Forward-Difference method is
consistent with local truncation error \(O(k+h^2)\), but consistency alone does
not guarantee useful numerical results. The method also needs
\(\lambda=k/h^2\leq 1/2\). In the stable case, the final-time profiles remain
smooth and the error decreases under refinement. In the unstable case, the
method amplifies high-frequency error modes instead of damping them, so the
computed solution can develop nonphysical oscillations even though the exact
heat equation smooths and decays the solution.

Backward Difference and Crank--Nicolson behave differently because the new time
level is included implicitly. Backward Difference is more diffusive, which can
make it slightly less accurate for the same mesh, but that diffusion is also why
it strongly damps unstable modes. Crank--Nicolson is less artificially
dissipative and has second-order temporal accuracy, so it usually gives a better
balance between accuracy and stability. The experiments show this tradeoff:
implicit methods require solving a linear system at every time step, but they
allow larger \(k\) without the Forward-Difference stability failure.

The Poisson manufactured-solution experiment is a better test of spatial
convergence than Problem A because \(u=\sin(\pi x)\sin(\pi y)\) is not exactly
represented by the five-point stencil. The observed errors decrease as the mesh
is refined, which matches the theoretical second-order behavior of the centered
spatial approximation. This experiment also shows why manufactured solutions
are useful: the source term is chosen so the exact solution is known, making it
possible to measure the numerical error directly.

The second additional experiment changes the heat equation initial condition to
\(\sin(2\pi x)\). This mode has a higher spatial frequency and decays like
\(e^{-4\pi^2t}\), much faster than the \(\sin(\pi x)\) mode. Numerically, this
means a coarse mesh has fewer points per wavelength, so spatial resolution
matters more. The experiment connects error propagation with the structure of
the PDE: the heat equation damps high-frequency physical modes, but a numerical
method can still introduce additional error if the grid or time step is too
coarse.

Overall, the experiments show the relationship between consistency, stability,
and convergence. The finite-difference formulas are consistent approximations,
but stability determines whether the accumulated error remains controlled. For
elliptic problems, the main computational issue is solving the sparse linear
system accurately and efficiently. For parabolic problems, the choice of time
step is just as important as the spatial mesh, especially for explicit methods.
Thus the best method depends not only on formal order, but also on stability
requirements and computational cost.

\section{Conclusion}

The experiments support the theoretical results for finite-difference PDE
methods. The five-point elliptic method is accurate and has a sparse matrix
structure that can be solved either directly or iteratively. For parabolic
problems, consistency alone is not enough: the Forward-Difference method also
requires the stability condition \(\lambda\leq 1/2\). Backward Difference and
Crank--Nicolson are more stable for large time steps, with Crank--Nicolson
offering higher temporal accuracy.

\end{document}
